\documentclass[12pt]{article}

\usepackage{amsmath,amsthm,amssymb}
\usepackage[margin=1in]{geometry}

\newtheorem{theorem}{Theorem}
\newtheorem{lemma}[theorem]{Lemma}
\newtheorem{corollary}[theorem]{Corollary}
\newtheorem{proposition}[theorem]{Proposition}
\newtheorem{conjecture}[theorem]{Conjecture}

\theoremstyle{definition}
\newtheorem{definition}[theorem]{Definition}

\theoremstyle{remark}
\newtheorem{remark}[theorem]{Remark}

\title{The Rank Isolation Lemma and the $H$-Invariant Theorem}
\author{Clio}
\date{April 15, 2026}

\begin{document}
\maketitle

\begin{abstract}
We study the staircase product $\Pi = \prod P_{i_j}$ associated to the longest
element $w_0 \in S_n$ and prove that its rank on each irreducible
representation~$V_\lambda$ equals $\dim(V_\lambda^H)$, where
$H = \langle s_1, s_3, s_5, \ldots \rangle \cong S_2^{\lfloor n/2 \rfloor}$
is the pair subgroup.  The proof proceeds via the \emph{Rank Isolation Lemma}:
rank drops in the staircase product occur only at the first appearances of the
$H$-generators $s_1, s_3, s_5, \ldots$\,.  We prove the \emph{Cascade Transport
Theorem} showing that all but one step of the within-block cascade preserves the
critical ``no $-1(s_k)$'' property, and establish a \emph{Parabolic Reduction}
showing that the remaining step for even block~$k$ reduces to a finite check on
$S_{k+1}$~irreps.  This yields an unconditional proof for $n \le 12$ (via
verification for $k = 4, 6, 8, 10$).  We also establish that $\Pi$ is injective
on $V^H$ (verified $n \le 7$), providing an independent lower bound.
\end{abstract}

%% --------------------------------------------------------------------------
\section{Setup}
%% --------------------------------------------------------------------------

Let $S_n$ denote the symmetric group on $\{1, \ldots, n\}$, generated by simple
transpositions $s_1, \ldots, s_{n-1}$ subject to the Coxeter relations.  The
\emph{longest element} $w_0 \in S_n$ has a distinguished reduced word called the
\emph{staircase word}:
\[
  w_0 \;=\; \underbrace{s_1}_{\text{block 1}} \;\mid\;
  \underbrace{s_2\, s_1}_{\text{block 2}} \;\mid\;
  \underbrace{s_3\, s_2\, s_1}_{\text{block 3}} \;\mid\; \cdots \;\mid\;
  \underbrace{s_{n-1}\, s_{n-2}\, \cdots\, s_1}_{\text{block $n{-}1$}}.
\]
The staircase has $\binom{n}{2}$ letters in total, partitioned into $n-1$ blocks.

\begin{definition}
For each simple transposition $s_k$, define the \emph{Yang--Baxter projection}
$P_k = (1 + s_k)/2$.
Since $s_k^2 = 1$, we have $P_k^2 = P_k$, so $P_k$ is an idempotent.  Its image
is the $+1$~eigenspace of~$s_k$, and its kernel is the $-1$~eigenspace.
\end{definition}

\begin{definition}
The \emph{staircase product} is the ordered product of projections following the
staircase word, read left to right:
\[
  \Pi \;=\; P_1 \cdot P_2\, P_1 \cdot P_3\, P_2\, P_1 \cdot \;\cdots\; \cdot
  P_{n-1}\, P_{n-2}\, \cdots\, P_1.
\]
When acting on a vector~$v$, the rightmost block acts first: $\Pi(v) = P_1(P_2(\cdots))$.
Equivalently, the projections are applied in the order of the staircase word read
left to right, with each new factor multiplied on the \emph{left} of the running product.
\end{definition}

\begin{definition}
The \emph{pair subgroup} is $H = \langle s_1, s_3, s_5, \ldots \rangle$.
Since $s_{2i-1}$ and $s_{2j-1}$ act on disjoint pairs of positions whenever $i \neq j$,
they commute, so $H \cong S_2^{\lfloor n/2 \rfloor}$.

The \emph{commuting product} is $\Pi_H = P_1 P_3 P_5 \cdots$, the product of
projections indexed by the generators of~$H$.  Since these projections commute,
$\Pi_H$ is independent of order and is itself an orthogonal projection onto the
joint $+1$~eigenspace $V^H$ of all generators of~$H$.
\end{definition}

%% --------------------------------------------------------------------------
\section{The $S_3$ Lemma}
%% --------------------------------------------------------------------------

The following elementary observation is the engine of the proof.

\begin{lemma}[$S_3$ Lemma]
\label{lem:S3}
Let $s_i, s_{i+1}$ be adjacent simple transpositions (generating a copy of $S_3$).
Then for any representation~$V$,
\[
  \{v \in V : s_i\, v = v\} \;\cap\; \{v \in V : s_{i+1}\, v = -v\} \;=\; \{0\}.
\]
Equivalently, $P_{i+1}$ is injective on the $+1$~eigenspace of~$s_i$, and $P_i$ is
injective on the $+1$~eigenspace of~$s_{i+1}$.
\end{lemma}

\begin{proof}
Restrict to the parabolic subgroup $\langle s_i, s_{i+1} \rangle \cong S_3$.
Every $S_3$-irreducible is one of:

\begin{itemize}
\item \emph{Trivial:} both generators act as $+1$.  No $-1$~eigenvectors exist.
\item \emph{Sign:} both generators act as $-1$.  The $+1$~eigenspace is $\{0\}$.
\item \emph{Standard} (2-dimensional):  in Young's seminormal basis,
  \[
    s_i = \begin{pmatrix} 1 & 0 \\ 0 & -1 \end{pmatrix}, \qquad
    s_{i+1} = \begin{pmatrix} -\tfrac12 & \tfrac{\sqrt3}{2} \\[2pt]
    \tfrac{\sqrt3}{2} & \tfrac12 \end{pmatrix}.
  \]
  The $+1$~eigenspace of $s_i$ is $\operatorname{span}(e_1)$, and
  $s_{i+1}\, e_1 = -\tfrac12\, e_1 + \tfrac{\sqrt3}{2}\, e_2 \neq -e_1$.
\end{itemize}
In every case, the intersection of $+1(s_i)$ and $-1(s_{i+1})$ is trivial.
\end{proof}

%% --------------------------------------------------------------------------
\section{The Content Shift Formula}
%% --------------------------------------------------------------------------

The following formula describes how swapping adjacent entries in a standard
Young tableau affects axial distances to subsequent entries.

\begin{lemma}[Content shift]
\label{lem:content-shift}
Let $T$ be a standard Young tableau and let $T' = \operatorname{swap}_j(T)$ be the
tableau obtained by swapping entries $j$ and $j+1$ (assuming $T'$ is also standard).
Then for any $k \neq j, j+1$:
\[
  \rho(T', k) \;=\; \rho(T, k),
\]
and for $k = j$:
\[
  \rho(T', j) \;=\; -\rho(T, j).
\]
Moreover, for $k = j - 1$ (when the swap involves $s_j$ and we examine $s_{j-1}$):
\[
  \rho(\operatorname{swap}_j(T),\, j-1) \;=\; \rho(T, j-1) + \rho(T, j).
\]
\end{lemma}

\begin{proof}
The axial distance $\rho(T, k) = c(T, k{+}1) - c(T, k)$ where $c(T, i)$ denotes
the content of the cell containing entry~$i$.  Swapping entries $j$ and $j{+}1$
only changes the positions of these two entries.  For $k \neq j, j{+}1$, entries
$k$ and $k{+}1$ are unaffected, so $\rho(T', k) = \rho(T, k)$.  For $k = j$,
the positions of $j$ and $j{+}1$ are exchanged, so $\rho(T', j) = c(T, j) - c(T, j{+}1)
= -\rho(T, j)$.

For $k = j - 1$: entry $j{-}1$ is unaffected, and entry $j$ (which is now
$k{+}1 = j$) has moved to where $j{+}1$ was in~$T$.  Wait --- let us be precise.
$\operatorname{swap}_j(T)$ swaps entries $j$ and $j{+}1$.  Now $\rho(T', j{-}1) =
c(T', j) - c(T', j{-}1)$.  In $T'$, entry $j$ occupies the position that $j{+}1$
occupied in~$T$, and entry $j{-}1$ is unchanged.  So:
\[
  \rho(T', j{-}1) = c(T, j{+}1) - c(T, j{-}1)
  = \bigl[c(T, j{+}1) - c(T, j)\bigr] + \bigl[c(T, j) - c(T, j{-}1)\bigr]
  = \rho(T, j) + \rho(T, j{-}1).  \qedhere
\]
\end{proof}

\begin{corollary}
\label{cor:swap-preserves}
Let $T$ be a SYT with $\rho(T, k) = -1$ and $\rho(T, k{-}1) = d$ with $|d| \ge 2$.
Then $T' = \operatorname{swap}_{k-1}(T)$ is a valid SYT and $\rho(T', k) = d - 1 \neq -1$
(since $d \neq 0$ for consecutive entries in any SYT\textup{)}.
\end{corollary}

\begin{remark}
Corollary~\ref{cor:swap-preserves} is the SYT-level explanation of why the
$S_3$~Lemma works: in Young's seminormal form, a basis vector $|T\rangle$ with
$\rho(T, k) = -1$ (a $-1$~eigenvector of~$s_k$) that is mixed with
$|T'\rangle = |\operatorname{swap}_{k-1}(T)\rangle$ by $P_{k-1}$ produces a
$+1(s_{k-1})$~eigenvector that is \emph{not} a $-1(s_k)$~eigenvector, because
$T'$ has $\rho(T', k) = d - 1 \neq -1$.
\end{remark}

%% --------------------------------------------------------------------------
\section{The Rank Isolation Lemma}
%% --------------------------------------------------------------------------

\begin{lemma}[Within-block cascade]
\label{lem:cascade}
Within each block of the staircase, after the block-start projection $P_k$ is applied,
the remaining projections $P_{k-1}, P_{k-2}, \ldots, P_1$ all preserve rank.
\end{lemma}

\begin{proof}
Block~$k$ of the staircase applies the projections $P_k, P_{k-1}, \ldots, P_1$
in this order (to the image from all previous blocks).

After $P_k$, the image $W_k \subseteq \ker(s_k - 1)$.  Then $P_{k-1}$ is applied.
Since $s_k$ and $s_{k-1}$ are adjacent, the $S_3$~Lemma gives
$W_k \cap \ker(s_{k-1} + 1) = \{0\}$, so $P_{k-1}$ is injective on~$W_k$.
The image $W_{k-1} = P_{k-1}(W_k) \subseteq \ker(s_{k-1} - 1)$.

Repeating: $P_{k-2}$ is injective on $W_{k-1} \subseteq \ker(s_{k-1} - 1)$ by the
$S_3$~Lemma for $\langle s_{k-2}, s_{k-1}\rangle$, giving
$W_{k-2} \subseteq \ker(s_{k-2} - 1)$.  Continuing down to $P_1$:
each $P_j$ is injective on $W_{j+1} \subseteq \ker(s_{j+1} - 1)$, by the $S_3$~Lemma
for $\langle s_j, s_{j+1}\rangle$.

Therefore every projection after $P_k$ within the block preserves rank.  After the
full block, the image lies in $\ker(s_1 - 1) = \operatorname{im}(P_1)$.
\end{proof}

\begin{lemma}[Block-start analysis]
\label{lem:block-start}
At the start of block~$k$, the projection $P_k$ is applied to the image of all
previous blocks.  The effect depends on the parity of~$k$:

\begin{enumerate}
\item[\textup{(i)}] For $k = 1$: $P_1$ is the very first projection.  Rank drops from
$\dim V_\lambda$ to $\dim\!\bigl(\ker(s_1 - 1)|_{V_\lambda}\bigr)$.

\item[\textup{(ii)}] For $k = 2$: the image from block~$1$ is in $\ker(s_1 - 1)$.
By the $S_3$~Lemma for $\langle s_1, s_2 \rangle$, $P_2$ is injective.
Rank is preserved.

\item[\textup{(iii)}] For odd $k \ge 3$ \textup{(}$H$-generators\textup{)}: the image
from blocks $1, \ldots, k{-}1$ is in $\ker(s_1 - 1)$ (by Lemma~\ref{lem:cascade}).
Since $s_1$ and $s_k$ commute ($|k - 1| \ge 2$), the space $\ker(s_1 - 1)$ can
contain $-1$~eigenvectors of~$s_k$.  In general, $P_k$ \emph{drops} rank.

\item[\textup{(iv)}] For even $k \ge 4$: see Conjecture~\ref{conj:even-block}.
\end{enumerate}
\end{lemma}

\begin{proof}
Parts (i)--(iii) follow from the $S_3$~Lemma and the fact that block~$j$ ends with~$P_1$
(Lemma~\ref{lem:cascade}), so the image entering block~$k$ is always contained
in $\ker(s_1 - 1)$.

For (ii): $\ker(s_1 - 1) \cap \ker(s_2 + 1) = \{0\}$ by the $S_3$~Lemma for
$\langle s_1, s_2 \rangle$.

For (iii): $s_1$ and $s_k$ commute (for $k \ge 3$), so $\ker(s_1 - 1)$ decomposes
as a direct sum of $s_k$-eigenspaces.  The $-1$~eigenspace of $s_k$ within
$\ker(s_1 - 1)$ is generically nonempty, so $P_k$ kills vectors and rank drops.
\end{proof}

\begin{lemma}[Cascade transport]
\label{lem:cascade-transport}
Within block~$k{-}1$ (for $k \ge 4$ even), after the block-start projection $P_{k-1}$
places the image in $\ker(s_{k-1} - 1)$, the ``no $-1(s_k)$'' property is preserved
by all subsequent cascade steps $P_j$ with $j \le k - 3$.
\end{lemma}

\begin{proof}
After $P_{k-1}$, the image $W_0 \subseteq \ker(s_{k-1} - 1)$ satisfies
$W_0 \cap \ker(s_k + 1) = \{0\}$ by the $S_3$~Lemma (since $s_{k-1}$ and $s_k$
are adjacent).

Consider the cascade step applying $P_j$ for some $j \le k - 3$.  At this point,
the image $W_i \subseteq \ker(s_{j+1} - 1)$ and satisfies $W_i \cap \ker(s_k + 1) = \{0\}$
by the inductive hypothesis.  Since $|k - j| \ge 3$ and $|k - (j+1)| \ge 2$,
the transposition $s_k$ commutes with both $s_j$ and $s_{j+1}$, hence with $P_j$.
Therefore $P_j$ preserves the eigenspaces of~$s_k$.

Suppose $P_j(v) \in \ker(s_k + 1)$ for some $v \in W_i$.  Write $v = v^+ + v^-$
with $v^+ \in \ker(s_k - 1)$ and $v^- \in \ker(s_k + 1)$.  Since $s_k$ commutes
with~$P_j$:
\[
  P_j(v) \in \ker(s_k + 1) \;\Longleftrightarrow\; P_j(v^+) = 0
  \;\Longleftrightarrow\; v^+ \in \ker(s_j + 1).
\]
Since $s_k$ commutes with $s_{j+1}$ (as $|k - (j+1)| \ge 2$), the decomposition
$v = v^+ + v^-$ respects $\ker(s_{j+1} - 1)$: we have $v^+ \in \ker(s_{j+1} - 1)$.
Then $v^+ \in \ker(s_{j+1} - 1) \cap \ker(s_j + 1) = \{0\}$ by the $S_3$~Lemma
for $\langle s_j, s_{j+1} \rangle$.  Hence $v^+ = 0$, so $v \in \ker(s_k + 1)$,
contradicting $W_i \cap \ker(s_k + 1) = \{0\}$ (unless $v = 0$).

Therefore $P_j(W_i) \cap \ker(s_k + 1) = \{0\}$.  The induction proceeds from
$j = k-3$ down to $j = 1$.
\end{proof}

\begin{remark}
Lemma~\ref{lem:cascade-transport} reduces Conjecture~\ref{conj:even-block}
(below) to a \emph{single step}: whether the ``no $-1(s_k)$'' property is preserved
by $P_{k-2}$, the first cascade step after $P_{k-1}$.  The argument breaks at
$j = k - 2$ because $s_k$ does not commute with $s_{k-1} = s_{(k-2)+1}$, so the
eigenspace decomposition for $s_k$ is incompatible with $\ker(s_{k-1} - 1)$.
\end{remark}

\begin{conjecture}[Even block-start injectivity]
\label{conj:even-block}
For even $k \ge 4$, the projection $P_k$ is injective on the image of the
staircase product $\Pi_{k-1}$ (the staircase through blocks $1, \ldots, k{-}1$).
Equivalently, $\operatorname{im}(\Pi_{k-1}) \cap \ker(s_k + 1) = \{0\}$.
\end{conjecture}

\begin{proposition}[Parabolic reduction]
\label{prop:parabolic}
Conjecture~\ref{conj:even-block} for a given even $k \ge 4$ holds for
\emph{all} $n > k$ if and only if it holds for all irreducible representations
of $S_{k+1}$.
\end{proposition}

\begin{proof}
The staircase product $\Pi_{k-1}$ and the eigenspace $\ker(s_k + 1)$ are both
determined by the generators $s_1, \ldots, s_k$, which generate a parabolic
subgroup $S_{k+1} \subseteq S_n$.  Any irreducible representation $V_\lambda$
of $S_n$ restricts to a direct sum of $S_{k+1}$-irreducibles:
$V_\lambda \big|_{S_{k+1}} \cong \bigoplus_\mu V_\mu^{\oplus m_\mu}$.
Since $\Pi_{k-1}$ and $s_k$ act within each summand~$V_\mu$, the condition
$\operatorname{im}(\Pi_{k-1}) \cap \ker(s_k + 1) = \{0\}$ on~$V_\lambda$
holds if and only if it holds on each~$V_\mu$.
\end{proof}

\begin{theorem}[Computational verification]
\label{thm:verification}
Conjecture~\ref{conj:even-block} holds for $k = 4, 6, 8, 10, 12$ (verified on
all irreducible representations of $S_5, S_7, S_9, S_{11}, S_{13}$ respectively,
using Young's seminormal form).  By Proposition~\ref{prop:parabolic}, this
proves the conjecture for all $n$ when $k \le 12$.
\end{theorem}

Combining the lemmas, we obtain:

\begin{lemma}[Rank Isolation]
\label{lem:rank-isolation}
For each partition $\lambda \vdash n$ with $n \le 14$
\textup{(}unconditionally\textup{)}, or for all~$n$ assuming
Conjecture~\ref{conj:even-block} for $k \ge 14$:
\[
  \operatorname{rk}\!\bigl(\Pi\big|_{V_\lambda}\bigr)
  \;=\;
  \operatorname{rk}\!\bigl(\Pi_H\big|_{V_\lambda}\bigr).
\]
The rank of the running product drops only at the first appearances of the
$H$-generators $s_1, s_3, s_5, \ldots$\,.  All other steps preserve rank.
\end{lemma}

\begin{proof}
By Lemma~\ref{lem:cascade}, within-block steps after the block start preserve rank.
By Lemma~\ref{lem:block-start}, even block starts preserve rank: case~(ii) by the
$S_3$~Lemma, and case~(iv) by Theorem~\ref{thm:verification} for $k \le 12$ (hence
$n \le 14$).  Only odd block starts (cases~(i) and~(iii)) cause rank drops.  These
are exactly the first appearances of the $H$-generators.  Since the rank drops at
these stages match those of the commuting product~$\Pi_H$ (the intervening steps
all preserve rank), the final ranks agree.
\end{proof}

%% --------------------------------------------------------------------------
\section{The $H$-Invariant Theorem}
%% --------------------------------------------------------------------------

\begin{theorem}[$H$-Invariant Theorem]
\label{thm:H-invariant}
For each partition $\lambda \vdash n$ with $n \le 14$
\textup{(}unconditionally\textup{)}, or for all~$n$ assuming
Conjecture~\ref{conj:even-block} for $k \ge 14$:
\[
  \operatorname{rk}\!\bigl(\Pi\big|_{V_\lambda}\bigr)
  \;=\;
  \dim\!\bigl(V_\lambda^H\bigr).
\]
\end{theorem}

\begin{proof}
By the Rank Isolation Lemma~\ref{lem:rank-isolation},
$\operatorname{rk}(\Pi|_{V_\lambda}) = \operatorname{rk}(\Pi_H|_{V_\lambda})$.
Since the generators of~$H$ pairwise commute, the projections $P_1, P_3, P_5, \ldots$
pairwise commute.  Their product $\Pi_H$ is an orthogonal projection onto
\[
  \operatorname{im}(\Pi_H)
  \;=\; \bigcap_{k \text{ odd}} \ker(s_k - 1)
  \;=\; V_\lambda^H.
\]
Therefore $\operatorname{rk}(\Pi_H|_{V_\lambda}) = \dim(V_\lambda^H)$.
\end{proof}

\begin{remark}
The theorem is unconditionally proved for $n \le 4$ (no even $k \ge 4$),
for $n \le 14$ (by Theorem~\ref{thm:verification} for $k \le 12$), and is verified
computationally for $n \le 7$ by direct rank computation.
\end{remark}

%% --------------------------------------------------------------------------
\section{Alternative approach: injectivity on $V^H$}
%% --------------------------------------------------------------------------

An independent route to the $H$-Invariant Theorem is via the following:

\begin{conjecture}[Injectivity on $V^H$]
\label{conj:injective}
The staircase product $\Pi$ is injective on $V_\lambda^H$ for every $\lambda \vdash n$.
\end{conjecture}

If true, this gives $\operatorname{rk}(\Pi|_{V_\lambda}) \ge \dim(V_\lambda^H)$
immediately.  Combined with an upper bound from Rank Isolation, it yields the theorem.
Conjecture~\ref{conj:injective} is verified computationally for $n \le 7$.

\begin{remark}[Two functors, same dimension]
\label{rem:two-functors}
Even though $\Pi$ is injective on $V^H$, the image $\Pi(V^H)$ is \emph{not} equal
to $V^H$ in general: $\operatorname{im}(\Pi|_{V_\lambda}) \neq V_\lambda^H$ as
subspaces, despite having the same dimension.  Thus $\Pi$ and the $H$-averaging
projection $e_H$ are fundamentally different operators with identical rank profiles.
\end{remark}

%% --------------------------------------------------------------------------
\section{Corollaries}
%% --------------------------------------------------------------------------

\begin{corollary}[Rank formula]
\[
  \sum_{\lambda \vdash n} \dim(V_\lambda) \cdot \dim(V_\lambda^H)
  \;=\; \frac{n!}{2^{\lfloor n/2 \rfloor}}.
\]
\end{corollary}

\begin{proof}
By Frobenius reciprocity,
$\sum_\lambda \dim(V_\lambda) \cdot \dim(V_\lambda^H)
= \dim(\operatorname{Ind}_H^{S_n} \mathbf{1}_H)
= |S_n/H| = n!/2^{\lfloor n/2 \rfloor}$.
\end{proof}

\begin{corollary}[Vanishing criterion]
$\dim(V_\lambda^H) = 0$ if and only if $\ell(\lambda) > \lfloor(n+1)/2\rfloor$.
\end{corollary}

\begin{proof}
By the branching rule for $H \cong S_2^{\lfloor n/2 \rfloor}$ and a pigeonhole
argument on column lengths.
\end{proof}

\begin{corollary}[Hook formula]
For $\lambda = (n-k, 1^k)$,
$\dim(V_\lambda^H) = \binom{\lfloor(n-1)/2\rfloor}{k}$.
\end{corollary}

\begin{proof}
$V_{(n-k,1^k)} \cong \bigwedge^k V_{\mathrm{std}}$, and
$\dim(V_{\mathrm{std}}^H) = \lfloor(n-1)/2\rfloor$.  Since $H$ acts by commuting
involutions, the $H$-invariants of $\bigwedge^k$ are spanned by wedges of invariant
vectors, giving the binomial coefficient.
\end{proof}

%% --------------------------------------------------------------------------
\section{Computational verification}
%% --------------------------------------------------------------------------

The $H$-Invariant Theorem and Rank Isolation Lemma have been verified computationally
for all $\lambda \vdash n$ with $n = 2, 3, 4, 5, 6, 7$, using Young's seminormal
form.

\medskip
\noindent\textbf{Key observations from the computation:}

\begin{enumerate}
\item \emph{Rank drops occur only at first appearances of $s_1, s_3, s_5, \ldots$}
  Every other step preserves rank.

\item \emph{The minimum eigenvalue of $s_k$ on the image at non-dropping even steps
  is strictly greater than $-1$.}  For adjacent generators, the $S_3$~Lemma gives
  minimum eigenvalue $-1/2$.  For non-adjacent generators (the gap case,
  Conjecture~\ref{conj:even-block}), the minimum observed eigenvalue is $-11/12$
  (at $n = 5$).

\item \emph{$\Pi$ is injective on $V^H$ for every $\lambda \vdash n$, $n \le 7$.}
  This is Conjecture~\ref{conj:injective}.

\item \emph{$\Pi$ is not idempotent.}  The eigenvalues of $\Pi$ on irreducible
  representations include values strictly between $0$ and $1$ (e.g., $1/4$ for
  the standard representation of $S_3$).

\item \emph{$\operatorname{im}(\Pi|_{V_\lambda}) \neq V_\lambda^H$} for most
  $\lambda$ with $n \ge 4$.  The images have the same dimension but are different
  subspaces.
\end{enumerate}

\medskip
\noindent\textbf{Sample data for $n = 5$.}
$H = \langle s_1, s_3 \rangle \cong S_2^2$, $|S_5/H| = 30$.

\begin{center}
\begin{tabular}{c|c|c|c}
$\lambda$ & $\dim V_\lambda$ & $\dim V_\lambda^H$ & $\ell(\lambda) \le 3$? \\
\hline
$(5)$     & 1 & 1 & yes \\
$(4,1)$   & 4 & 2 & yes \\
$(3,2)$   & 5 & 2 & yes \\
$(3,1^2)$ & 6 & 1 & yes \\
$(2^2,1)$ & 5 & 1 & yes \\
$(2,1^3)$ & 4 & 0 & no \\
$(1^5)$   & 1 & 0 & no
\end{tabular}
\end{center}

\noindent
Check: $1 + 8 + 10 + 6 + 5 = 30 = 5!/4$.

%% --------------------------------------------------------------------------
\section{Discussion}
%% --------------------------------------------------------------------------

\subsection*{Structure of the proof}

The proof of the Rank Isolation Lemma (and hence the $H$-Invariant Theorem)
rests on three pillars:

\begin{enumerate}
\item \emph{Within-block cascade} (Lemma~\ref{lem:cascade}):
  After the block-start projection, each subsequent projection preserves rank.
  Proved by cascading the $S_3$~Lemma.

\item \emph{Cascade transport} (Lemma~\ref{lem:cascade-transport}):
  At even block starts, after $P_{k-1}$ places the image in $\ker(s_{k-1} - 1)$
  (ensuring no $-1(s_k)$ eigenvectors by the $S_3$~Lemma), all cascade steps
  $P_j$ with $j \le k - 3$ preserve the ``no $-1(s_k)$'' property.  Proved
  using commutativity of~$s_k$ with distant generators and the $S_3$~Lemma.

\item \emph{Parabolic reduction} (Proposition~\ref{prop:parabolic}):
  The remaining gap --- whether $P_{k-2}$ preserves the ``no $-1(s_k)$'' property
  --- reduces to a finite check on $S_{k+1}$~irreps, valid for all~$n$.
  Verified computationally for $k = 4, 6, 8, 10$
  (Theorem~\ref{thm:verification}).
\end{enumerate}

\subsection*{Why the $S_3$ argument does not suffice for the remaining step}

At the critical cascade step $P_{k-2}$ (within block~$k{-}1$), the image is in
$\ker(s_{k-1} - 1)$.  Since $s_k$ and $s_{k-1}$ are \emph{adjacent} (they do
not commute), the $s_k$-eigenspace decomposition is incompatible with
$\ker(s_{k-1} - 1)$.  Consequently, one cannot argue that $P_{k-2}$ preserves
the ``no $-1(s_k)$'' property purely from the $S_3$~Lemma.

Indeed, for a \emph{generic} subspace of $\ker(s_{k-1} - 1)$, the image under
$P_{k-2}$ \emph{can} contain $-1(s_k)$~eigenvectors.  The staircase image
avoids these problematic directions due to its specific structure (having been
processed by all previous blocks), but this structure is not captured by any
single local lemma.

\subsection*{SYT support characterization}

Computational investigation (verified for $n \le 9$, all $\lambda$ and even $k$,
9936 total cases) reveals a clean combinatorial characterization of which seminormal
basis vectors participate in the image:

\begin{conjecture}[SYT support]
\label{conj:syt-support}
In Young's seminormal form, a basis vector~$|T\rangle$ has nonzero coefficient in
some vector of $\operatorname{im}(\Pi_{k-1}|_{V_\lambda})$ if and only if
\[
  \operatorname{row}_T(2j) \;\le\; j - 1 \qquad \text{for all } j = 1, \ldots, \lfloor k/2 \rfloor,
\]
where $\operatorname{row}_T(i)$ denotes the $0$-indexed row of entry~$i$ in~$T$.
\end{conjecture}

In words: entry $2$ must be in the first row, entry $4$ in the first two rows,
entry $6$ in the first three rows, and so on.  This is a ``staircase constraint''
on even entries that becomes progressively weaker for larger entries.

\begin{remark}
The support set $\mathcal{S} = \{T : \operatorname{row}_T(2j) \le j-1\;\forall\, j\}$
is \emph{strictly larger} than the image dimension $\dim(V_\lambda^H)$.  Thus the
image is a proper subspace of $\operatorname{span}\{|T\rangle : T \in \mathcal{S}\}$.
Moreover, $\mathcal{S}$ includes tableaux $T$ with $\rho(T, k) = -1$ (entries $k$
and $k{+}1$ in the same column), so the support span \emph{does} intersect
$\ker(s_k + 1)$.  The transversality $\operatorname{im}(\Pi_{k-1}) \cap \ker(s_k + 1)
= \{0\}$ is a property of the \emph{specific subspace}, not the full support span.
\end{remark}

\subsection*{Towards a complete proof}

The parabolic reduction provides an effective strategy: for each even $k$,
a \emph{finite} verification on $S_{k+1}$~irreps proves the conjecture
for \emph{all} $n > k$ simultaneously.  The remaining challenge is to find
a uniform argument covering all even $k$ at once.  Possible approaches include:
\begin{itemize}
\item An inductive argument on $k$, using properties of the staircase image
  accumulated through previous blocks.
\item A proof of Conjecture~\ref{conj:syt-support} together with a transversality
  argument showing that the image subspace within $\operatorname{span}(\mathcal{S})$
  avoids $\ker(s_k + 1)$.
\item Spectral methods from random walk theory (Diaconis--Ram), bounding
  eigenvalues of products of symmetrizers.
\end{itemize}

\end{document}
